\subsection{Appendix: an incomplete argument}\label{sy-0300-ai}

This file is nested, so its section renders one level down.

\begin{lemma}[Incomplete]\label{sy-000C-ai}
Every gadget has even order.
\end{lemma}
\begin{proof}
\uses{sy-0008-ai}
Orbits have size two. \incomplete{Show that there are no fixed points using Definition~\ref{sy-0008-ai}.}
\end{proof}

\begin{lemma}[No proof yet]\label{sy-000D-ai}
Every widget map preserves orbit sizes.
\end{lemma}

\begin{remark}
An unlabelled remark: gadgets never have odd order.
\end{remark}

\begin{definition}[Half-gadget]\label{sy-000E-ai}
A widget with exactly one fixed point.
\end{definition}
\begin{proof}
Definitions need no proof; this one has a stray proof anyway.
\end{proof}

% !LOOM priority: high

\begin{proof}[Proof of Lemma~\ref{sy-000D-ai} and Lemma~\ref{sy-0007-ai}]
A proof naming two results attaches to the first.
\end{proof}


% A missing \input for the missing-include diagnostic. LaTeX would abort on it, so it sits inside \iffalse, which LaTeX skips and the scanner does not interpret.
\iffalse
\fi



\begin{mainthm}\label{sy-000B-ai}
\uses{sy-0003-ai, sy-0006-ai}
Every finite widget of odd order has a fixed point, and the fixed locus is functorial.
\end{mainthm}
\begin{proof}
Combine Theorem~\ref{sy-0003-ai} and Proposition~\ref{sy-0006-ai}.
\end{proof}


\bibliographystyle{amsplain}
\bibliography{refs}

